\documentclass[12pt,a4paper]{report}

% Packages
\usepackage[utf8]{inputenc}
\usepackage[english]{babel}
\usepackage{graphicx}
\usepackage{amsmath}
\usepackage{amssymb}
\usepackage{hyperref}
\usepackage{cite}
\usepackage{listings}
\usepackage{xcolor}
\usepackage{geometry}
\usepackage{fancyhdr}
\usepackage{setspace}

% Page geometry
\geometry{
    left=3cm,
    right=2.5cm,
    top=2.5cm,
    bottom=2.5cm
}

% Header and footer
\pagestyle{fancy}
\fancyhf{}
\rhead{\thepage}
\lhead{\leftmark}

% Code listing settings
\lstset{
    basicstyle=\ttfamily\small,
    breaklines=true,
    frame=single,
    numbers=left,
    numberstyle=\tiny\color{gray},
    keywordstyle=\color{blue},
    commentstyle=\color{green!60!black},
    stringstyle=\color{orange},
    showstringspaces=false
}

% Hyperref setup
\hypersetup{
    colorlinks=true,
    linkcolor=blue,
    filecolor=magenta,
    urlcolor=cyan,
    citecolor=blue,
}

% Document information
\title{Internship Assignment Optimization System}
\author{Your Name}
\date{\today}

\begin{document}

% Title page
\maketitle

% Management Summary (non-technical, decision-level)
\chapter*{Management Summary}
\addcontentsline{toc}{chapter}{Management Summary}

\section*{Client and Context}
% The university unit responsible for organizing teaching internships
% Challenge: coordinating hundreds of students across multiple schools

\section*{Core Challenges}
% • Large number of students requiring placements
% • Multiple internship types (PDP, ZSP, SFP) with different subjects
% • Limited teacher availability and capacity
% • Distance, zone, and subject matching constraints
% • Budget constraints limiting number of active placements

\section*{Solution Delivered}
% A web-based software system that:
% • Automates teacher and student assignment to internship placements
% • Optimizes assignments across multiple planning phases
% • Provides configurable constraints and budget controls
% • Offers interactive assignment management and validation

\section*{Outcomes}
% • High-quality, feasible assignment plans
% • Significant reduction in manual planning effort
% • Transparent, auditable decision-making
% • Flexible configuration for changing requirements

% Table of contents
\tableofcontents
\newpage

% Chapter 1: Introduction (motivation & context)
\chapter{Introduction}

\section{Teaching Internships in Teacher Education}
% • Role of practical experience in teacher training programs
% • Different internship types serving different pedagogical goals
% • Importance of quality supervision by experienced teachers

\section{Scale and Complexity of Internship Coordination}
% • Hundreds of students per semester
% • Dozens of schools and supervising teachers
% • Multiple subject areas and school types (GS, MS)
% • Geographic distribution and accessibility considerations

\section{Limitations of Manual Planning}
% • Time-intensive process
% • Difficulty balancing competing constraints
% • Prone to errors and inconsistencies
% • Lack of optimization and transparency
% • Challenge of reoptimization when circumstances change

\section{Course Context}
% • Real client: university internship coordination unit
% • Multiple student teams working on related aspects
% • Agile, sprint-based development over multiple semesters
% • Collaboration with stakeholders and teaching assistants

\section{Team Contribution}
% • Your specific role and deliverables
% • Integration with work from other teams
% • Evolution of the system across development cycles

% Chapter 2: Requirements Overview
\chapter{Requirements Overview}

\section{System Vision}
% • Support annual internship planning cycle
% • Optimize assignments under real-world constraints
% • Prioritize solution quality over strict feasibility guarantees
% • Provide transparency and configurability

\section{Stakeholders}
\subsection{Primary: Internship Coordination Unit}
% • Needs efficient planning tools
% • Requires audit trails and justification for decisions

\subsection{Teaching Students}
% • Need fair, quality placements
% • Prefer locations close to home
% • Require subject-appropriate supervision

\subsection{Schools and Supervising Teachers}
% • Have capacity and availability constraints
% • Bring subject expertise and preferences
% • Constrained by geographic zones and accessibility

\section{Core Problem Definition}

\subsection{Input Data}
% Students:
% • Internship type requirements (PDP_I, PDP_II, ZSP, SFP)
% • Subject preferences and specializations
% • Home addresses for distance calculations

% Schools:
% • Type (GS - elementary, MS - middle school)
% • Location, zone, and public transport accessibility
% • Available supervising teachers

% Teachers:
% • Subject expertise and qualifications
% • Capacity and availability preferences
% • Maximum workload constraints

\subsection{Constraints}
% Hard constraints (must be satisfied):
% • Budget: total number of active internship slots
% • Teacher workload: valid combinations (1, 2, or 4 internships)
% • Subject matching: teachers must be qualified for assigned subjects
% • Zone feasibility: appropriate geographic zones per internship type
% • Course pinning: certain subjects fixed, others flexible

% Soft constraints (preferences to optimize):
% • Minimize student travel distance (PDP internships)
% • Balance teacher workload (prefer 2 internships per teacher)
% • Diversity: mix of internship types at each school
% • Preserve teacher assignments across semesters
% • Zone preferences: optimal zones per internship type

\subsection{Objectives}
% • Produce high-quality assignment plans
% • Maximize satisfaction of soft constraints
% • Acknowledge: not guaranteed feasibility due to conflicting constraints
% • System is an optimization system rather than a pure feasibility solver

% Chapter 3: Architecture Overview
\chapter{Architecture Overview}

\section{System Architecture}
% • Backend-centered design
% • RESTful API for frontend integration
% • Database for persistent storage
% • Optimization engine integration

\section{Technology Stack}

\subsection{Backend}
% • Spring Boot framework
% • OptaPlanner constraint satisfaction solver
% • PostgreSQL database
% • Maven build system

\subsection{Frontend}
% • React with TypeScript
% • Vite build tool
% • Component-based architecture

\section{Separation of Concerns}
% • Data preparation and validation layer
% • Assignment optimization logic
% • User-facing management interfaces
% • Audit and reporting subsystems

\section{Core Data Model}
% • Students and their internship requirements
% • Teachers with qualifications and availability
% • Schools with capacity and characteristics
% • Internship slots (planning entities)
% • Assignments (relationships)

\section{Multi-Phase Assignment Strategy}
% The assignment logic is structured into multiple stages to reflect 
% real-world planning constraints:
% • Preparation: transform requirements into demand slots
% • Teacher assignment: allocate teachers and activate slots
% • Student assignment: match students to activated slots
% Rationale: separates capacity planning from individual assignment

% Chapter 4: System Implementation
\chapter{System Implementation}

\section{Preparation Phase (Phase 0)}

\subsection{Demand Slot Generation}
% • Transform student requirements into demand slots
% • One slot per internship type per school type
% • Aggregate student demand by type and location

\subsection{Data Enrichment}
% • Address geocoding for distance calculations
% • Subject preference mapping
% • School capacity and teacher availability validation

\subsection{Rationale}
% • Decouples demand modeling from individual student assignment
% • Enables capacity planning before granular assignments
% • Simplifies teacher assignment optimization

\section{Phase 1: Teacher Assignment and Slot Activation}

\subsection{Optimization Goal}
% • Assign supervising teachers to internship slots
% • Activate optimal subset of slots within budget
% • Maximize coverage while respecting capacity constraints

\subsection{Planning Variables}
% For each slot:
% • assignedTeacher (Teacher or null)
% • active (boolean)
% • course (for flexible course assignments like ZSP)

\subsection{Key Constraints}
% Hard constraints:
% • Budget: total active slots ≤ budget limit
% • Teacher workload: 0, 1, 2, or 4 internships per teacher
% • Teacher diversity: if 2+ internships, all different types
% • Subject qualification: teachers must be qualified for assigned courses
% • Zone feasibility: appropriate zones per internship type
% • Course pinning: PDP has no course, SFP course is fixed

% Soft constraints:
% • Prefer 2 internships per teacher
% • Minimize PDP distance (teachers to student demand centroids)
% • Diversity: varied internship types at each school
% • Preserve teacher assignments from previous semester

\subsection{Distance Optimization for PDP}
% Two mechanisms:
% 1. Bidirectional matching: each student finds closest teacher, 
%    each teacher finds closest student
% 2. Centroid distance: teachers positioned toward student demand center
% Ensures geographic spread and accessibility

\subsection{Trade-offs}
% • Balance coverage vs. budget
% • Trade off teacher preference vs. student demand
% • Optimize for quality, may not achieve strict feasibility

\section{Phase 2: Student Assignment}

\subsection{Optimization Goal}
% • Assign individual students to activated slots
% • Minimize total distance for PDP students
% • Respect subject preferences

\subsection{Planning Variables}
% For each student:
% • assignedSlot (PlannedInternship or null)

\subsection{Key Constraints}
% Hard constraints:
% • Each student assigned to exactly one slot
% • Slot must match student's internship type requirement
% • Slot must be active (from Phase 1)
% • Subject matching for ZSP/SFP

% Soft constraints:
% • Minimize distance for PDP students to assigned school
% • Prefer slots matching student subject preferences

\subsection{Rationale}
% • Separation from Phase 1 reduces complexity
% • Teacher availability already determined
% • Focus purely on student satisfaction
% • Enables what-if scenarios without reoptimizing teachers

\section{Optimization Engine}

\subsection{OptaPlanner Integration}
% • Constraint satisfaction solver
% • Hard/soft score calculation
% • Local search and metaheuristics
% • Termination conditions and time limits

\subsection{Score Calculator Implementation}
% • Manual score calculator to avoid caching issues
% • Pure Java constraint evaluation
% • Efficient incremental scoring
% • Detailed logging for debugging

\subsection{Solver Configuration}
% • Local search algorithms (Tabu Search, Late Acceptance)
% • Time budget per optimization run
% • Score weight tuning

% Chapter 5: Project Process
\chapter{Project Process}

\section{Development Methodology}
% • Agile/Scrum process
% • Sprint duration and planning
% • Stakeholder engagement

\section{Evolution of the System}

\subsection{Sprint Overview}
% • Sprint 1-2: Initial data model and basic assignment
% • Sprint 3-4: Distance optimization and zone constraints
% • Sprint 5-6: Multi-phase optimization and preservation
% • Sprint 7+: Refinement and testing

\subsection{Key Milestones}
% • When Phase 0/1/2 separation was introduced
% • Distance optimization added after client feedback
% • Preservation constraint for semester continuity
% • Budget and workload constraint refinements

\section{Team Collaboration}
% • Coordination with other teams
% • Integration of shared components
% • Code review and quality practices

\section{Quality Metrics}
% • GitLab activity and commit patterns
% • TeamScale code quality scores
% • Test coverage
% • Defect tracking

\section{Challenges and Solutions}
% • Technical challenges (OptaPlanner caching, performance)
% • Requirement changes and adaptations
% • Balancing optimization quality vs. runtime

% Chapter 6: Results and Evaluation
\chapter{Results and Evaluation}

\section{Test Scenarios}
% • Small-scale validation datasets
% • Real-world data from previous semesters
% • Edge cases and stress tests

\section{Performance Analysis}

\subsection{Phase 1 Solver Performance}
% • Runtime vs. problem size
% • Score improvement over time
% • Feasibility rate

\subsection{Phase 2 Solver Performance}
% • Student assignment success rate
% • Distance metrics
% • Optimization convergence

\section{Solution Quality}

\subsection{Constraint Satisfaction}
% • Hard constraint violations (should be zero)
% • Soft constraint scores
% • Comparison to manual assignments (if available)

\subsection{Stakeholder Feedback}
% • Client satisfaction
% • Usability observations
% • Practical deployment considerations

\section{Limitations}
% • Cases where feasibility could not be achieved
% • Trade-offs in multi-objective optimization
% • Computational constraints for large-scale problems

% Chapter 7: Conclusion
\chapter{Conclusion}

\section{Summary}
% • Recap of problem and solution
% • Key achievements
% • Technical contributions

\section{Reflection on Optimization Goals}
% • Were optimization objectives met?
% • How well did multi-phase approach work?
% • Balance between automation and manual control

\section{Lessons Learned}
% • Technical insights
% • Process learnings
% • Collaboration experiences

\section{Future Work}

\subsection{Technical Enhancements}
% • Better multi-objective optimization
% • Machine learning for constraint tuning
% • Real-time reoptimization

\subsection{Feature Extensions}
% • Scenario comparison and what-if analysis
% • Advanced reporting and visualization
% • Integration with university information systems

\subsection{Process Improvements}
% • Automated testing and validation
% • Performance optimization
% • Scalability for larger datasets

% Bibliography
\bibliographystyle{plain}
\bibliography{references}

% Appendices
\appendix
\chapter{Additional Materials}
\section{Code Samples}
% Add important code snippets

\section{Configuration Files}
% Add configuration examples

\end{document}
